\documentclass[12pt, a4paper]{article}

% --- PAQUETES ESENCIALES ---
\usepackage[utf8]{inputenc}
\usepackage[T1]{fontenc}
\usepackage[spanish]{babel}
\usepackage{amsmath}
\usepackage{graphicx}
\usepackage{booktabs}
\usepackage{geometry}
\usepackage{hyperref}
\usepackage{fancyhdr}
\usepackage[table]{xcolor}
\usepackage{listings}

% --- ETIQUETA 'LISTING' A 'LISTADO' ---
\addto{\captionsspanish}{\renewcommand{\lstlistingname}{Listado}}
\addto{\captionsspanish}{\renewcommand{\lstlistlistingname}{Listado de Códigos}}

% --- ESTILO DE CÓDIGO ---
\definecolor{GrisClaro}{rgb}{0.95, 0.95, 0.95}

\lstset{
    backgroundcolor=\color{GrisClaro},
    basicstyle=\ttfamily\small,
    breaklines=true,
    showstringspaces=false,
    numbers=left,
    numberstyle=\tiny\color{gray},
    frame=single,
    framesep=5pt,
    rulesepcolor=\color{black},
    keywordstyle=\color{blue}\bfseries,
    commentstyle=\color{green!60!black}\itshape,
    stringstyle=\color{orange},
    tabsize=2,
}

% Lenguajes que listings no trae nativamente
\lstdefinelanguage{Kotlin}{
    keywords={package, import, fun, override, val, var, when, class, return, super, private, public, internal, object, this, if, else, true, false, null, in, is, as},
    keywordstyle=\color{blue}\bfseries,
    ndkeywords={String, Int, Boolean, Unit, Bundle, MethodChannel, FlutterActivity, FlutterEngine, WindowManager, Location},
    ndkeywordstyle=\color{purple}\bfseries,
    sensitive=true,
    comment=[l]{//},
    morecomment=[s]{/*}{*/},
    morestring=[b]",
    morestring=[b]',
}

\lstdefinelanguage{Dart}{
    keywords={import, class, abstract, static, final, const, void, return, if, else, async, await, Future, bool, String, int, double, true, false, null, this, super, extends, new, try, catch, on, in, is, as, override, enum, switch, case, default},
    keywordstyle=\color{blue}\bfseries,
    ndkeywords={Widget, State, StatefulWidget, StatelessWidget, BuildContext, MaterialApp, MethodChannel, Platform, PlatformException, Geolocator, Position, LocationPermission, LocationAccuracy, LocationSettings, Duration, IntegrityResult, IntegrityStatus, MockLocationGuard},
    ndkeywordstyle=\color{purple}\bfseries,
    sensitive=true,
    comment=[l]{//},
    morecomment=[s]{/*}{*/},
    morestring=[b]",
    morestring=[b]',
}

\lstdefinelanguage{XML}{
    morestring=[b]",
    morestring=[s]{>}{<},
    morecomment=[s]{<?}{?>},
    morecomment=[s]{<!--}{-->},
    stringstyle=\color{black},
    identifierstyle=\color{purple},
    keywordstyle=\color{blue}\bfseries,
    morekeywords={xmlns,version,encoding,manifest,application,activity,uses-permission,android:name,android:label,intent-filter,action,category,meta-data,queries,intent,data}
}

% --- MÁRGENES ---
\geometry{
    a4paper,
    total={170mm,257mm},
    left=25mm,
    right=25mm,
    top=30mm,
    bottom=30mm,
}

% --- ENCABEZADO Y PIE ---
\pagestyle{fancy}
\fancyhf{}
\renewcommand{\headrulewidth}{0pt}
\rfoot{\thepage}
\lfoot{Seguridad de la Información / Miguel Ángel Gutiérrez Gómez}

% --- DATOS DEL DOCUMENTO ---
\newcommand{\nombreuniversidad}{Universidad Politécnica de Chiapas}
\newcommand{\nombrecarrera}{Ingeniería en Software}
\newcommand{\nombremateria}{Seguridad de la Información}
\newcommand{\nombreprofesor}{MC. José Alonso Macías Montoya}
\newcommand{\corte}{Corte 1}
\newcommand{\actividad}{Práctica - Móvil 2 \\ Detección de Ubicación Falsa \\ (Bloqueo de Fake GPS)}
\newcommand{\nombres}{Miguel Ángel Gutiérrez Gómez}
\newcommand{\matricula}{233382}
\newcommand{\fechaentrega}{\today}
\newcommand{\ciudad}{Suchiapa, Chiapas}

\begin{document}

% --- PORTADA ---
\begin{titlepage}
    \centering
    \vspace*{1cm}
    {\Huge\bfseries \nombreuniversidad \par}
    \vspace{0.5cm}
    {\Large \nombrecarrera \par}

    \vspace{2cm}
    \rule{\textwidth}{1.5pt}
    \vspace{0.4cm}
    {\Large\bfseries \actividad \par}
    \vspace{0.4cm}
    \rule{\textwidth}{1.5pt}

    \vspace{3cm}

    \begin{flushleft}
        \centering
        \textbf{Materia:} \nombremateria \\
        \textbf{Profesor:} \nombreprofesor \\
        \textbf{Corte:} \corte
    \end{flushleft}

    \vspace{0.5cm}

    \begin{flushleft}
        \centering
        \textbf{Alumno:} \nombres \\
        \textbf{Matrícula:} \matricula \\
    \end{flushleft}

    \vfill

    \vspace{1cm}
    \textbf{\ciudad, a \fechaentrega}

\end{titlepage}

\clearpage

% --- ÍNDICE ---
\tableofcontents
\clearpage

% --- CUERPO ---

\section{Introducción}
\label{sec:introduccion}
Las aplicaciones de Fake GPS son herramientas que permiten al usuario
\textit{inyectar} coordenadas arbitrarias en el sistema operativo, de
manera que cualquier aplicación que consulte la ubicación del dispositivo
recibirá un valor falso en lugar del real obtenido por GPS, Wi-Fi o red
celular. Este tipo de manipulación es un vector de fraude habitual en
aplicaciones de reparto a domicilio, registro de asistencia laboral,
juegos basados en geolocalización, banca digital con verificación por
ubicación, transporte y servicios de auditoría en campo.

El objetivo de esta práctica es evitar que la aplicación móvil
desarrollada se ejecute en un dispositivo en el que se encuentre activo
un Fake GPS, deteniendo el acceso a las pantallas posteriores
(incluido el login) y mostrando una pantalla informativa de bloqueo. La
detección se hace mediante las APIs nativas de Android para reportar
ubicaciones simuladas (\texttt{isMock()} a partir de API 31 y
\texttt{isFromMockProvider()} desde API 18), expuestas en Flutter a
través del paquete \texttt{geolocator}.

\section{Recursos Utilizados}
\label{sec:recursos}
A continuación se listan las herramientas y tecnologías empleadas:

\begin{itemize}
    \item \textbf{SDK:} Flutter 3.44.0 (Dart 3.12.0).
    \item \textbf{Lenguajes:} Dart (capa Flutter) y Kotlin (capa nativa Android, mantenida igual que en la práctica anterior).
    \item \textbf{Plataforma objetivo:} Android.
    \item \textbf{Paquete:} \texttt{geolocator} (manejo de permisos y posiciones, expone \texttt{Position.isMocked}).
    \item \textbf{APIs nativas Android:} \texttt{Location.isMock()} (API 31+) y \texttt{Location.isFromMockProvider()} (API 18+).
    \item \textbf{Permisos:} \texttt{ACCESS\_FINE\_LOCATION} y \texttt{ACCESS\_COARSE\_LOCATION}.
    \item \textbf{Editor:} Visual Studio Code.
    \item \textbf{Entorno de pruebas:} Dispositivo Android físico con app de Fake GPS instalada y emulador de Android Studio.
\end{itemize}

\section{Desarrollo}
\label{sec:desarrollo}

\subsection{Cómo funciona la detección}
Cuando una aplicación de Fake GPS está activa y registrada como
\textit{Mock Location App} en las Opciones de Desarrollador de Android,
las ubicaciones que entrega el sistema a través de
\texttt{LocationManager} o de Google Play Services traen una bandera
interna que marca la posición como simulada. Esta bandera puede leerse:

\begin{itemize}
    \item En \textbf{API 31 (Android 12) o superior} con
    \texttt{Location.isMock()}.
    \item En \textbf{API 18 a 30} con
    \texttt{Location.isFromMockProvider()} (deprecada en API 31, pero
    aún funcional como compatibilidad).
\end{itemize}

Cualquiera de las dos retorna \texttt{true} si la coordenada fue
inyectada por un proveedor de prueba (test provider) o por una
aplicación con permiso \texttt{ACCESS\_MOCK\_LOCATION}. El paquete
\texttt{geolocator} encapsula ambas y las expone en Dart como
\texttt{Position.isMocked}, que es la lectura que se utiliza en esta
práctica.

\subsection{Arquitectura de la solución}
La solución se compone de cuatro piezas:

\begin{enumerate}
    \item \textbf{Permisos en el manifest}
    (\texttt{AndroidManifest.xml}): declaración de
    \texttt{ACCESS\_FINE\_LOCATION} y \texttt{ACCESS\_COARSE\_LOCATION}.
    \item \textbf{Servicio de verificación}
    (\texttt{mock\_location\_guard.dart}): orquesta la solicitud de
    permisos, la obtención de la posición y la lectura de
    \texttt{isMocked}. Devuelve un objeto \texttt{IntegrityResult} con
    uno de seis estados posibles (\texttt{ok}, \texttt{mockDetected},
    \texttt{locationServiceDisabled}, \texttt{permissionDenied},
    \texttt{permissionDeniedForever}, \texttt{unknownError}).
    \item \textbf{Pantalla de integridad / splash}
    (\texttt{integrity\_gate\_screen.dart}): es la pantalla inicial de
    la app. Muestra el logo más un indicador de progreso y, tras el
    primer frame, invoca al servicio. Según el resultado redirige a
    \texttt{/login} o a \texttt{/blocked}.
    \item \textbf{Pantalla de bloqueo}
    (\texttt{blocked\_screen.dart}): comunica al usuario el motivo
    del bloqueo, lista los pasos para resolverlo y ofrece los botones
    \textit{Reintentar} y \textit{Cerrar aplicación}.
\end{enumerate}

\subsection{Flujo de ejecución}
\begin{enumerate}
    \item Al lanzar la app se abre \texttt{IntegrityGateScreen}.
    \item El servicio comprueba que la plataforma sea Android. En iOS y
    desktop devuelve \texttt{ok} (no existe API equivalente para
    detectar mock).
    \item Verifica que el servicio de ubicación esté encendido.
    \item Solicita el permiso de ubicación si aún no se ha otorgado.
    \item Obtiene la última posición conocida; si no hay caché, solicita
    una posición actual con precisión mínima y tiempo límite de 8 s.
    \item Lee \texttt{position.isMocked}.
    \item Si es \texttt{true}, devuelve \texttt{mockDetected} y la
    aplicación navega a \texttt{BlockedScreen}. Si es \texttt{false},
    navega a \texttt{LoginScreen}.
\end{enumerate}

\subsection{Implementación}

\begin{lstlisting}[language=XML, caption={Permisos de ubicación en AndroidManifest.xml}, label={lst:manifest}]
<manifest xmlns:android="http://schemas.android.com/apk/res/android">
    <!-- Permisos requeridos para detectar uso de Fake GPS / mock locations -->
    <uses-permission android:name="android.permission.ACCESS_FINE_LOCATION"/>
    <uses-permission android:name="android.permission.ACCESS_COARSE_LOCATION"/>

    <application android:label="scanbuild" ...>
        <!-- resto del manifest sin cambios -->
    </application>
</manifest>
\end{lstlisting}

\begin{lstlisting}[caption={Dependencia agregada en pubspec.yaml}, label={lst:pubspec}]
dependencies:
  flutter:
    sdk: flutter
  cupertino_icons: ^1.0.8
  geolocator: ^13.0.1
\end{lstlisting}

\begin{lstlisting}[language=Dart, caption={Servicio de verificación de integridad}, label={lst:guard}]
import 'dart:io';
import 'package:geolocator/geolocator.dart';

enum IntegrityStatus {
  ok,
  mockDetected,
  locationServiceDisabled,
  permissionDenied,
  permissionDeniedForever,
  unsupportedPlatform,
  unknownError,
}

class IntegrityResult {
  final IntegrityStatus status;
  final String? detail;
  const IntegrityResult(this.status, [this.detail]);

  bool get isBlocking =>
      status == IntegrityStatus.mockDetected ||
      status == IntegrityStatus.locationServiceDisabled ||
      status == IntegrityStatus.permissionDenied ||
      status == IntegrityStatus.permissionDeniedForever;
}

class MockLocationGuard {
  static Future<IntegrityResult> check() async {
    if (!Platform.isAndroid) {
      return const IntegrityResult(IntegrityStatus.ok,
          'Plataforma sin deteccion de mock locations.');
    }

    try {
      final serviceEnabled = await Geolocator.isLocationServiceEnabled();
      if (!serviceEnabled) {
        return const IntegrityResult(
          IntegrityStatus.locationServiceDisabled,
          'El servicio de ubicacion esta apagado. Activalo para continuar.',
        );
      }

      var permission = await Geolocator.checkPermission();
      if (permission == LocationPermission.denied) {
        permission = await Geolocator.requestPermission();
      }
      if (permission == LocationPermission.denied) {
        return const IntegrityResult(
          IntegrityStatus.permissionDenied,
          'La app necesita permiso de ubicacion para verificar la integridad.',
        );
      }
      if (permission == LocationPermission.deniedForever) {
        return const IntegrityResult(
          IntegrityStatus.permissionDeniedForever,
          'El permiso de ubicacion esta bloqueado permanentemente.',
        );
      }

      Position? position = await Geolocator.getLastKnownPosition();
      position ??= await Geolocator.getCurrentPosition(
        locationSettings: const LocationSettings(
          accuracy: LocationAccuracy.lowest,
          timeLimit: Duration(seconds: 8),
        ),
      );

      if (position.isMocked) {
        return const IntegrityResult(
          IntegrityStatus.mockDetected,
          'Se detecto una ubicacion falsa (Fake GPS / mock location).',
        );
      }

      return const IntegrityResult(IntegrityStatus.ok);
    } catch (e) {
      return IntegrityResult(IntegrityStatus.unknownError, e.toString());
    }
  }
}
\end{lstlisting}

\begin{lstlisting}[language=Dart, caption={Pantalla de integridad (splash) que evalúa al iniciar}, label={lst:gate}]
class _IntegrityGateScreenState extends State<IntegrityGateScreen> {
  @override
  void initState() {
    super.initState();
    WidgetsBinding.instance.addPostFrameCallback((_) => _evaluate());
  }

  Future<void> _evaluate() async {
    final result = await MockLocationGuard.check();
    if (!mounted) return;

    if (result.isBlocking) {
      Navigator.of(context).pushReplacementNamed(
        '/blocked',
        arguments: result,
      );
    } else {
      Navigator.of(context).pushReplacementNamed('/login');
    }
  }
}
\end{lstlisting}

\begin{lstlisting}[language=Dart, caption={Registro de rutas con el gate como punto de entrada}, label={lst:main}]
return MaterialApp(
  title: 'VisionPrice',
  debugShowCheckedModeBanner: false,
  theme: AppTheme.light(),
  initialRoute: '/',
  routes: {
    '/'        : (_) => const IntegrityGateScreen(),
    '/blocked' : (_) => const BlockedScreen(),
    '/login'   : (_) => const LoginScreen(),
    '/home'    : (_) => const HomeShell(),
    '/forgot'  : (_) => const ForgotPasswordScreen(),
    '/register': (_) => const RegisterScreen(),
  },
);
\end{lstlisting}

\subsection{Evidencia visual}
La Figura \ref{fig:gate} muestra la pantalla de verificación de
integridad mientras se consulta el estado del proveedor de ubicación.
La Figura \ref{fig:bloqueo-fake} muestra la pantalla de bloqueo
desplegada cuando se detecta una aplicación de Fake GPS activa.
La Figura \ref{fig:login-ok} muestra el acceso permitido al login
cuando la verificación se completa sin detectar ubicación falsa.

\begin{figure}[h]
    \centering
    % \includegraphics[width=0.45\textwidth]{gate.png}
    \fbox{\parbox[c][6cm][c]{0.45\textwidth}{\centering \textit{Splash de verificación de integridad}}}
    \caption{Pantalla de verificación al iniciar la aplicación.}
    \label{fig:gate}
\end{figure}

\begin{figure}[h]
    \centering
    % \includegraphics[width=0.45\textwidth]{bloqueo_fake.png}
    \fbox{\parbox[c][6cm][c]{0.45\textwidth}{\centering \textit{Pantalla de acceso bloqueado por Fake GPS}}}
    \caption{Acceso bloqueado al detectar una aplicación de Fake GPS activa.}
    \label{fig:bloqueo-fake}
\end{figure}

\begin{figure}[h]
    \centering
    % \includegraphics[width=0.45\textwidth]{login_ok.png}
    \fbox{\parbox[c][6cm][c]{0.45\textwidth}{\centering \textit{Login mostrado tras verificación OK}}}
    \caption{Login mostrado cuando la verificación termina sin detectar mock locations.}
    \label{fig:login-ok}
\end{figure}

\section{Resultados y Conclusiones}
\label{sec:resultados}

\subsection{Resultados}
Se realizaron tres pruebas representativas sobre el dispositivo Android:

\begin{itemize}
    \item \textbf{Caso A — Dispositivo con Fake GPS activo:} se instaló
    una aplicación de Fake GPS, se eligió en
    \textit{Ajustes \(\rightarrow\) Opciones de Desarrollador \(\rightarrow\)
    Seleccionar app de ubicación simulada} y se activó una posición
    arbitraria. Al lanzar la aplicación, el guard reportó
    \texttt{mockDetected} y se mostró la \texttt{BlockedScreen} con el
    mensaje de bloqueo y los pasos de remediación. No se permitió
    avanzar al login.
    \item \textbf{Caso B — Dispositivo sin Fake GPS:} con todas las apps
    de ubicación simulada desinstaladas o desactivadas, el guard
    devolvió \texttt{ok} y la aplicación pasó directamente a la
    pantalla de login.
    \item \textbf{Caso C — Servicio de ubicación apagado:} con el GPS
    del sistema apagado, el guard devolvió
    \texttt{locationServiceDisabled} y se mostró la pantalla de bloqueo
    indicando que es necesario activar la ubicación para verificar la
    integridad del dispositivo.
\end{itemize}

\subsection{Conclusiones}
\begin{itemize}
    \item Se cumplió el objetivo de la práctica: la aplicación no se
    ejecuta (no permite avanzar al login) en un dispositivo donde se
    está utilizando un Fake GPS.
    \item Colocar la verificación en una pantalla de \textit{gate} antes
    del login garantiza que la lógica de autenticación nunca quede
    expuesta en un contexto comprometido.
    \item El uso del paquete \texttt{geolocator} abstrae las diferencias
    entre \texttt{isMock()} y \texttt{isFromMockProvider()} y deja un
    único punto de lectura (\texttt{Position.isMocked}), lo que simplifica
    el mantenimiento.
    \item Limitaciones honestas:
    \begin{itemize}
        \item En iOS no existe una API pública para detectar
        ubicaciones simuladas; la defensa en esa plataforma debe
        complementarse con validaciones server-side.
        \item En el emulador de Android Studio, las posiciones
        inyectadas mediante los controles del emulador son técnicamente
        mock, por lo que para pruebas funcionales del login conviene
        usar un dispositivo físico o aceptar el bloqueo como parte
        del comportamiento esperado.
        \item Un dispositivo con root podría falsificar la bandera
        \texttt{isMock}; para mayor robustez se recomienda
        complementar con detección de root, validaciones del lado
        del servidor (coordenadas vs. IP, patrones de movimiento) y
        atestación de integridad como Google Play Integrity API.
    \end{itemize}
\end{itemize}

\clearpage

\section{Anexos}
\label{sec:anexos}

\subsection{Estructura de archivos relevantes}

\begin{lstlisting}[language={}, caption={Archivos clave de la solución anti Fake GPS}]
scanbuild/
  android/app/src/main/
    AndroidManifest.xml          <-- Permisos de ubicacion
  lib/
    main.dart                    <-- Rutas (initialRoute = '/')
    services/
      mock_location_guard.dart   <-- Verificacion de isMocked
    screens/
      integrity_gate_screen.dart <-- Splash de verificacion
      blocked_screen.dart        <-- Pantalla de bloqueo
      login_screen.dart          <-- Login (alcanzable solo si OK)
  pubspec.yaml                   <-- Dependencia geolocator
\end{lstlisting}

\subsection{Procedimiento de prueba}
\begin{enumerate}
    \item Instalar en el dispositivo una aplicación de Fake GPS
    (por ejemplo \textit{Fake GPS Location}).
    \item Activar las Opciones de Desarrollador en Android
    (\textit{Ajustes \(\rightarrow\) Acerca del teléfono \(\rightarrow\)} pulsar
    7 veces sobre el número de compilación).
    \item Dentro de Opciones de Desarrollador entrar a
    \textit{Seleccionar app de ubicación simulada} y elegir la app de
    Fake GPS instalada.
    \item Abrir la app de Fake GPS, fijar una coordenada arbitraria y
    pulsar \textit{Start}.
    \item Lanzar la aplicación VisionPrice. Se debe observar la
    pantalla de bloqueo con el mensaje \textit{``Acceso bloqueado''}.
    \item Detener la app de Fake GPS o seleccionar \textit{Ninguna} en
    la opción de ubicación simulada y volver a abrir la aplicación.
    Debe observarse el flujo normal hasta el login.
\end{enumerate}

\subsection{Evidencia adicional}
A continuación se reserva espacio para capturas complementarias de la
app de Fake GPS configurada en las Opciones de Desarrollador y del
dispositivo mostrando el bloqueo.

\begin{figure}[h]
    \centering
    % \includegraphics[width=0.45\textwidth]{dev_options.png}
    \fbox{\parbox[c][6cm][c]{0.45\textwidth}{\centering \textit{App de Fake GPS configurada en Opciones de Desarrollador}}}
    \caption{Configuración de la app de ubicación simulada en Opciones de Desarrollador.}
    \label{fig:devopts}
\end{figure}

\end{document}
